% DeepPass68: finite-dimensional Gaussian and Grassmann identities.

多变量高斯公式直接由实对称正定二次型定义；一维积分只是其最简单的分量检验。对实对称正定矩阵 $A$，定义带源积分
\begin{equation}
\mathcal I_A[J]\equiv
\int_{\mathbb R^n}\dd^nq\,
\exp\!\left[-\frac12q^TAq+J^Tq\right].
\label{eq:euclidean-gaussian-definition-dp68}
\end{equation}
完成平方后得到
\begin{equation}
\boxed{
\mathcal I_A[J]
=(2\pi)^{n/2}(\det A)^{-1/2}
\exp\!\left[\frac12J^TA^{-1}J\right]}
\label{eq:euclidean-gaussian-source-dp21}
\end{equation}
以及归一化二点平均
\begin{equation}
\boxed{
\langle q_iq_j\rangle_A=(A^{-1})_{ij}.}
\label{eq:gaussian-covariance-dp21}
\end{equation}
这里 $\langle\cdots\rangle_A$ 表示以 $\ee^{-q^TAq/2}$ 为权重并除以零源积分后的平均。\eqnrefs{eq:gaussian-covariance-dp21} 体现了二次理论的基本结构：二次核 $A$ 控制作用，而二点响应由它的逆矩阵 $A^{-1}$ 控制。

费米自由度需要另一种积分。Grassmann 变量满足 $\theta_i\theta_j=-\theta_j\theta_i$，所以特别有 $\theta_i^2=0$。Berezin 积分不是面积积分，而是抽取最高次 Grassmann 系数的线性运算；一维约定固定为
\begin{equation}
\int\dd\theta\,1=0,
\qquad
\int\dd\theta\,\theta=1.
\label{eq:berezin-definition-dp68}
\end{equation}
对 $n$ 对独立变量 $\theta_i,\bar\theta_i$，固定测度次序后定义
\begin{equation}
\mathcal Z_G[\xi,\bar\xi]\equiv
\int\prod_{i=1}^{n}\dd\bar\theta_i\dd\theta_i\,
\exp\!\left[-\bar\theta M\theta+\bar\xi\theta+\bar\theta\xi\right],
\label{eq:grassmann-source-definition-dp68}
\end{equation}
其中 $M$ 可逆，而 $\xi,\bar\xi$ 也是 Grassmann 源。采用与\eqnrefs{eq:grassmann-source-definition-dp68} 一致的固定测度约定，零源积分为
\begin{equation}
\boxed{
\mathcal Z_G[0,0]=\det M.}
\label{eq:grassmann-gaussian-determinant-dp68}
\end{equation}
加入源以后，平移变量给出
\begin{equation}
\boxed{
\frac{\mathcal Z_G[\xi,\bar\xi]}{\mathcal Z_G[0,0]}
=\exp\!\left(\bar\xi M^{-1}\xi\right).}
\label{eq:grassmann-source-gaussian-dp68}
\end{equation}
因此，若对 $\bar\xi_i$ 取左导数、对 $\xi_j$ 取右导数，则零源二点收缩为
\begin{equation}
\boxed{
\left.
\frac{1}{\mathcal Z_G[0,0]}
\frac{\overrightarrow\partial\mathcal Z_G}{\partial\bar\xi_i}
\frac{\overleftarrow\partial}{\partial\xi_j}
\right|_{\xi=\bar\xi=0}
=(M^{-1})_{ij}.}
\label{eq:grassmann-inverse-kernel-dp68}
\end{equation}
这正是费米子自由理论中“二次算符的逆成为传播核”的有限维原型；闭合费米子圈的额外负号则来自 Wick 收缩中的费米置换宇称，与高斯积分行列式是不同的代数来源。

最后还需要知道 Grassmann 测度怎样随线性变量变换改变。对
\begin{equation}
\theta_i'=A_{ij}\theta_j,
\label{eq:berezin-linear-transform-dp68}
\end{equation}
固定同一测度次序后有
\begin{equation}
\boxed{
\dd\theta_n'\cdots\dd\theta_1'
=(\det A)^{-1}\,
\dd\theta_n\cdots\dd\theta_1.}
\label{eq:berezin-jacobian-dp21}
\end{equation}
它与普通变量的 Jacobian 方向相反，来源就是\eqnrefs{eq:berezin-definition-dp68} 的系数抽取规则。
